%=======================================================================
\section{Introduction}
\label{sec:intro}
%=======================================================================
% Paragraphs 1--4 are carried over from the IHTC-18 paper with light edits.
% Paragraphs 5--7 are new and state what this paper adds.

The expansion of variable renewable energy (VRE) sources is worsening
temporal mismatches between electricity supply and demand. Grid operators
increasingly face renewable curtailment during periods of peak generation and
the underutilisation of assets owing to inadequate long-duration energy
storage (LDES) \cite{denholm2021,iea2020}. Ensuring grid reliability while
meeting decarbonisation targets therefore requires firm, dispatchable,
low-emission technologies alongside VRE, including enhanced geothermal
systems (EGS) and subsurface storage technologies such as compressed air
energy storage (CAES), gravitational energy storage (GES) and thermal energy
storage (TES). Although EGS holds long-term promise, high upfront costs,
drilling risk and long development timelines constrain near-term deployment
\cite{doe2006,doe2019}.

Repurposing idle and mature oil and gas wells for LDES offers a near-term
pathway that leverages existing infrastructure while reducing environmental
liabilities. Borehole-based TES stores surplus electricity during
over-generation and releases it during peak demand \cite{bartling2025}.
Reusing existing wells avoids the most capital-intensive element of
geothermal development, new drilling, and thereby lowers the levelised cost
of storage \cite{santos2022,higgins2024}. Idle and orphaned wells are also a
significant source of fugitive methane \cite{kang2014,williams2021,Santos2025};
integrating them into closed-loop storage mitigates that risk and defers
abandonment costs.

Several authors have proposed legacy wells for energy storage. Zhang
\textit{et al.}\ \cite{Zhang2025} evaluated geothermal-assisted CAES in
plugged wells and reported round-trip efficiencies of about \SI{61}{\percent}
for an output of about \SI{37}{\mega\watt} over a \SI{4}{\hour} discharge, a
system requiring of the order of 100 wells. Jaquez \cite{Jaquez2021} proposed
plugged wells as sealed shafts for GES, and drivetrain studies of that concept
reported round-trip efficiencies of up to \SI{85.9}{\percent}
\cite{Sundeep}.

In a preceding conference paper \cite{KoIHTC2026} the authors proposed a
power-to-heat-to-power (P2H2P) Carnot battery in which a plugged,
hydraulically isolated well is filled with a phase-change material (PCM) and
equipped with finned hairpin tubes (the ``borehole battery'', BB). A
high-temperature heat pump (HTHP) charges the BB; an organic Rankine cycle
(ORC) discharges it. The principal novelty was an axial \emph{cascade} of PCM
layers of decreasing melting temperature, arranged so that the secondary
fluid meets material whose melting point tracks its own temperature glide.
PCM-enhanced borehole heat exchangers have been studied for ground-source heat
pumps \cite{Qi2016,Zhou2025}, but not, to the authors' knowledge, as the
storage element of a Carnot battery with a graded melting temperature.

The borehole model of Ref.~\cite{KoIHTC2026} advanced a melt-layer thickness
from a closed-form quasi-steady Stefan solution with a prescribed wall
temperature, evaluated a single charge from a fully solid store, and sized the
field by root-finding on the number of wells until the charged energy matched
the heat-pump output. Three features of that formulation limit its use for a
storage system operated on a repeating daily schedule, and they motivate the
present work.

\begin{enumerate}[leftmargin=1.8em,itemsep=2pt]
\item \textbf{State variable.} A melt thickness (or melted area) encodes only
  how much PCM has changed phase. Once a segment has melted completely the
  resistance network continues to report a heat flow, and the model must
  either melt material that does not exist or discard the heat. Neither is
  acceptable in a cascaded store, where the segments nearest the inlet of each
  layer saturate early.
\item \textbf{Single cycle.} A store operated daily does not start each
  charge from a cold, fully solid state. The operating point is the cyclic
  steady state (CSS) reached after many cycles, and first-cycle results are
  not representative of it.
\item \textbf{Plant--store coupling.} The ORC evaporating temperature was
  set by a fixed approach below the hot end of the water stream. For a
  working fluid that absorbs most of its heat isothermally against a water
  stream with a glide of several tens of kelvin, the binding constraint is an
  interior pinch, and the hot-end rule overstates the cycle efficiency.
\end{enumerate}

This paper reformulates the model to remove these limitations. The
contributions are: (i)~a lumped \emph{enthalpy} state for each exchanger
segment, in the tradition of the enthalpy method for phase-change problems
\cite{Voller1987,Alexiades1993}, which represents subcooled, two-phase and
superheated states with one invertible constitutive relation; (ii)~the
derivation of the segment conductance from control-volume balances, which
yields an effectiveness--NTU closure rather than the product of the overall
coefficient and the perimeter; (iii)~a time integrator that is exact within
each branch of the constitutive relation and therefore unconditionally stable
on the logarithmic grids natural to diffusion-limited problems; (iv)~a
direction-dependent conduction path for melting and freezing, and a
mean-to-surface resistance for single-phase segments; (v)~irreversible cycle models in
which the isentropic efficiencies act on the state points, a pinch-constrained
coupling of the store to the ORC evaporator and an energy-consistent closure
of the HTHP regenerators; (vi)~a reformulation of sizing as a simulation
to cyclic steady state, which is necessary because the storage efficiency of
an adiabatic store is identically unity at CSS; and (vii)~a component-wise
exergy audit, verified machine by machine against independent boundary
balances, in which the low-temperature geothermal source of the heat pump is
treated as a finite, priced resource rather than as a free reservoir. The model is described in
Section~\ref{sec:model}, its numerical implementation in
Section~\ref{sec:implementation}, and results for a
\SI{1}{\mega\watt}$_\mathrm{e}$, \SI{10}{\hour} pilot configuration in
Section~\ref{sec:results}.
